% Appendix for insertion into a larger paper.
% Most preparatory material is copied from wilson.tex; original NB comments retained.
\section{The derived Ravenel--Wilson Hopf ring}\label{app:rw-derived}
\providecommand{\gpt}[1]{\todo[color=cyan!25,linecolor=cyan!60!black,bordercolor=cyan!60!black]{\textcolor{teal!80!black}{#1}}}
\providecommand{\Free}{\operatorname{Free}}
\providecommand{\Ga}{{\mathbb G}_a}
\providecommand{\RW}{\mathrm{RW}}
\providecommand{\dRW}{\scr{RW}}
\providecommand{\Sym}{\mathrm{Sym}}
\providecommand{\LSym}{\operatorname{LSym}}
\providecommand{\Vect}{\mathrm{Vect}}
\providecommand{\hotimes}{\odot}
\providecommand{\htimes}{\otimes}
\providecommand{\hH}{\mathfrak H}
\providecommand{\frob}{\mathrm f}
\providecommand{\ver}{\mathrm v}
\providecommand{\wittprojection}{\mathrm \pi}
\providecommand{\HCAlg}{\mathrm{HCAlg}}
\providecommand{\Hopf}{\mathrm{Hopf}}
\providecommand{\MU}{\mathrm{MU}}
\providecommand{\Set}{\mathrm{Set}}
\providecommand{\Ab}{\mathrm{Ab}}
\providecommand{\cn}{\mathrm{cn}}
\providecommand{\cl}{\pi}
\providecommand{\A}{\mathbb A}
\providecommand{\cof}{\operatorname{cof}}
\providecommand{\comp}{\wedge}
\providecommand{\wequi}{\simeq}
\providecommand{\ph}{\mathord-}
\providecommand{\fr}{\mathrm{fr}}
\providecommand{\Hom}{\operatorname{Hom}}
\providecommand{\CAlg}{\mathrm{CAlg}}
\providecommand{\cCAlg}{\mathrm{cCAlg}}
\providecommand{\bCAlg}{\mathrm{bCAlg}}
\providecommand{\1}{\mathbbm{1}}

In this section we introduce the \emph{derived Ravenel--Wilson Hopf ring} and show that it coincides with the ordinary Ravenel--Wilson Hopf ring of \cite{ravenel-wilson}.
The material has been extracted from our forthcoming article \cite{bachmann-wilson} in order to make the current manuscript more self-contained.

We fix a prime $p$ and work over $\F_p$.
We write $L=L_{(p)}$ for the $p$-local Lazard ring.  We write $\Hopf_{\F_p}$ for the category of evenly graded bicommutative Hopf algebras over $\F_p$. We write $H\hotimes K$ for the tensor product of Hopf algebras, so maps $H\hotimes K\to T$ correspond to bilinear coalgebra maps $H\htimes K\to T$. A \emph{Hopf ring} is a commutative ring object for $\hotimes$, with a separate external grading.
The Hopf ring associated with a graded commutative ring $A$ is denoted $\F_p[A]$.

If $H$ is a Hopf algebra, $\Sym(H)$ denotes the free commutative Hopf ring on $H$ placed in external degree one; its component of external degree $n$ is $H^{\hotimes n}/\Sigma_n$. Similarly $\Lambda(H)$ denotes the exterior algebra for $\hotimes$, obtained by killing the image of the diagonal. For a graded coefficient ring $A$, we write $\Sym_A(H)=\F_p[A]\hotimes\Sym(H)$.

We regard a simplicial Hopf algebra as a simplicial object in the additive category of graded Hopf algebras. Its homology, written $\hH_i$, is the homology of the normalized chain complex under Dold--Kan. The same convention applies degreewise to simplicial Hopf rings.

\subsection{Frobenius, Verschiebung, and Dieudonn\'e theory}\label{app:rw-witt}

% COPIED from wilson.tex, lines 823--861; source comments retained.
For each $n>0$ we have the functor $\phi_n^*: \Vect_{\F_p}^\N \to \Vect_{\F_p}^\N$ defined by $(\phi_n^*V)_{ni} = V_i$ and $(\phi_n^*V)_j = 0$ for $j$ not a multiple of $n$.
The functor $\phi_n^*$ admits an evident symmetric monoidal structure and hence extends canonically to $\cCAlg(\Vect_{\F_p}^\N)$, and so on.
Moreover, $\phi_n^*$ is symmetric monoidal for $\hotimes$.

Given $A \in \CAlg(\Vect_{\F_p})$ we have the Frobenius endomorphism $F: A \to A, a \mapsto a^p$.
This is in fact a symmetric monoidal endotransformation of the identity functor.
When working with graded algebras, the Frobenius changes grading, and so for $A \in \CAlg(\Vect_{\F_p}^\N)$ takes the form $F: \phi_p^* A \to A$.
Being symmetric monoidal, this upgrades to \[ F:\phi_p^* \to \id \in \Fun(\Hopf_{\F_p}, \Hopf_{\F_p}). \]

We shall also need the \emph{Verschiebung} map, which arises in a dual fashion.
For a (graded) vector space $C$, denote by $C^{tp}$ the quotient of $(C^{\otimes p})^{C_p}$ by the subspace generated by vectors of the form $\bigoplus_{\sigma \in C_p} \sigma v$.
Then there is a canonical equivalence of (oplax symmetric monoidal) endofunctors $(\ph)^{tp} \wequi \phi_p^*$.\NB{ref?}
Consequently for a graded coalgebra we obtain a map \[ V: C \xrightarrow{\Delta^p} (C^{\otimes p})^{h\Sigma_p} \to (C^{\otimes p})^{tp} \wequi \phi_p^* C \] called the \emph{Verschiebung}.
Again this is a symmetric monoidal endotransformation and hence induces \[ V: \id \to \phi_p^*  \Fun(\Hopf_{\F_p}, \Hopf_{\F_p}). \]

We recall the following important facts.
\begin{lemma} \label{lemm:FV-basics}
\begin{enumerate}
\item We have $FV = [p] = VF$.
\item Given $H, K \in \Hopf_{\F_p}$ we have $V_{H \hotimes K} = V_H \hotimes V_K$ and the following diagram commutes
\begin{equation*}
\begin{CD}
\phi_p^*H \hotimes \phi_p^*K @>{\wequi}>> \phi_p^*(H \hotimes K) \\
@A{\id_H \hotimes V_K}AA                   @V{F_{H \hotimes K}}VV \\
\phi_p^* H \hotimes K @>{F_H \hotimes \id_K}>> H \hotimes K.
\end{CD}
\end{equation*}
\end{enumerate}
\end{lemma}
\begin{proof}
(1) That $FV = [p]$ is essentially the point of the definition, and follows from the fact that $F$ vanishes on elements of $(H^{\otimes p})^{C_p}$ which are \emph{induced}, i.e., in the image of $(H^{\otimes p})_{C_p}$.
Naturality of $V$ applied to the map $F: \phi_p^* H \to H$ shows that $V_H F_H = F_{\phi_p^* H} V_{\phi_p^* H} = [p]_{\phi_p^* H}$, as needed.

(2) It suffices to check the first equality on free Hopf algebras, where it reduces to the analogous claim for coalgebras and the usual tensor product, which is clear (by construction).
The second statement follows from the distributivity of $\hotimes$ over $\htimes$.
(In slightly more detail, given $x \in H_{i/p}$  and $y \in K_j$ we can express $x^p \hotimes y$ in terms of the iterated diagonal $\Delta^p(y)$, writing out this formula precisely yields the desired commutative square.)
\end{proof}

% COPIED from wilson.tex, lines 867--882; source comments retained.
Given a Hopf algebra $H$, applying indecomposables to Frobenius and Verschiebung yields maps \[ QF: \phi_p^* QH \to QH \quad\text{and}\quad QV: QH \to \phi_p^* QH. \]
Clearly $QF = 0$ (i.e., $a^p$ is always decomposable).
This is not so for $QV$ and so the graded vector space $QH$ acquires an additional structure, namely the Verschiebung endomorphism.
Let us write $\Vect_{\F_p}^{\N,V}$ for the category of such vector spaces, i.e., graded vector spaces $X$ together with a map $V: X \to \phi_p^* X$.

\begin{proposition} \label{prop:indecomp-charp}
The above construction upgrades to a symmetric monoidal functor \[ Q: \Hopf_{\F_p} \to \Vect_{\F_p}^{\N,V} \] which restricts to an equivalence \[ \Hopf_{\F_p}^{\cn,F=0} \wequi (\Vect_{\F_p}^{\N,V})^\cn. \]
Here by $\Hopf_{\F_p}^{\cn,F=0} \subset \Hopf_{\F_p}^{\cn}$ we denote the subcategory of those Hopf algebras whose Frobenius endomorphism vanishes, i.e., $H$ such that $F: \phi_p^* H \to H$ factors through $\F_p$.
\end{proposition}
\begin{proof}
A symmetric monoidal structure is inherited from the universal property of $\hotimes$ and Lemma \ref{lemm:FV-basics}(2).
The fact that the restricted functor is an equivalence is part of Dieudonné theory: there is an equivalence $D$ from $\Hopf_{\F_p}^{\cn}$ to a category of graded modules over $\Z_p$ with an action by $F$ and $V$ such that $FV = VF = p$.
Moreover there is a canonical short exact sequence $D\phi_p^* H \xrightarrow{DF} DH \to QH \to 0$.
See, e.g., \cite[Theorem 4.7, Lemma 4.4]{goerss-dieudonne}.
The result follows.
\end{proof}

% COPIED from wilson.tex, lines 884--889; source comments retained.
\subsubsection{Additive group} \label{subsec:Ga}
An important Hopf algebra for us will be $\Ga$, the Hopf algebra corresponding to the additive group over $\F_p$.
As a ring it is given by \[ \Ga = \F_p[a_1, a_2, a_3, \dots]/(a_i a_j = {i+j \choose j} a_{i+j}), \text{ where } |a_i| = i. \]
The coalgebra structure is determined by \[ \Delta(a_i) = \sum_{i = j + k} a_j \otimes a_k. \]
The underlying coalgebra of $\Ga$ is the base change to $\F_p$ of the $L$-coalgebra $\A^1$ defined in Definition~\ref{def:A1}. We use the same symbol for its base change whenever the coefficient ring is clear.

% COPIED from wilson.tex, lines 898--908; source comments retained.
\begin{example} \label{ex:Ga-hotimes-Ga}
We now find $a_i^p = 0$\NB{details?} and \[ \Ga \wequi \F_p[a_{(0)}, a_{(1)}, \dots]/a_{(i)}^p, \text{ where } a_{(i)} = a_{p^i}. \]
It follows that $F=0$.
One may check that $V: \Ga \to \phi_p^* \Ga$ sends $a_{(0)}$ to $0$, whereas for $i>0$ it sends $a_{(i)}$ to an element which might be denoted $\phi_p^*a_{(i-1)}$.
In particular $V$ is surjective and we tend to denote the generators of $\phi_p^* \Ga$ by their names coming from $\Ga$, i.e., we tend to think of the element $\phi_p^*a_{(i-1)}$ as $Va_{(i)}$ and even denote it by just $a_{(i)}$.
We thus have \[ Q\phi_{p^n}^*\Ga \wequi \F_p\{a_{(n)}, a_{(n+1)}, \dots\} \text{ with } Va_{(i)} = a_{(i-1)} \text{ and } a_{(i)} = 0 \text{ for } i<n. \]
Using the equivalence between connected Hopf algebras and graded vector spaces with Verschiebung (Proposition \ref{prop:indecomp-charp}) we see that, for example, \begin{gather*} \Ga \hotimes \Ga \wequi \F_p[a_{(i)} \circ a_{(j)}, i,j \ge 0]/(a_{(i)} \circ a_{(j)})^p = 0 \text{ with } V(a_{(i)} \circ a_{(j)}) = a_{(i-1)} \circ a_{(j-1)} \\ \text{ where } a_{(i)} \circ a_{(j)} = 0 \text{ whenever $i<0$ or $j<0$}. \end{gather*}
Equivalently \[ \Ga \hotimes \Ga \wequi \bigotimes_{i > 0} \phi_{p^i+1}^*\Ga \otimes \phi_2^*\Ga \otimes \bigotimes_{j > 0} \phi_{p^j+1}^*\Ga, \] with the copies of $\Ga$ corresponding respectively to $a_{(i)} \circ a_{(0)}$, $a_{(0)} \circ a_{(0)}$ and $a_{(0)} \circ a_{(j)}$.
Many similar computations are possible and straightforward, e.g., \[ \Sym^2 \Ga \wequi \F_p[a_{(i)} \circ a_{(j)}, i,j \ge 0]/(a_{(i)} \circ a_{(j)})^p = 0,a_{(i)} \circ a_{(j)} = a_{(j)} \circ a_{(i)}. \]
\end{example}

% COPIED from wilson.tex, lines 910--946; source comments retained.
\subsubsection{Hopf algebras with a lift of Verschiebung} \label{subsec:lift-of-v}
Following Goerss, by a \emph{Hopf algebra with a lift of Verschiebung} we mean a pair of
\begin{itemize}
\item an object $H \in \HCAlg(\Mod_{\Z_p}^{\heart\N})$, and
\item a Hopf algebra map $V: H \to \phi_p^* H$ such that $V \otimes \F_p = V_{H \otimes \F_p}$.
\end{itemize}
We denote the category of these objects by $\HCAlg_{\Z_p}^V$.
It carries a symmetric monoidal structure, i.e., a version of $\hotimes$.
Indeed given two Hopf algebras over $\Z_p$ $H_1, H_2$, with lifts of Verschiebung $V_1, V_2$, the map $V_1 \hotimes V_2$ lifts the Verschiebung on $H_1 \hotimes H_2$ (Lemma \ref{lemm:FV-basics}(2)).

Let us write $\Mod_{\Z_p}^{\N,V}$ for the category of graded $\Z_p$-modules together with an action by $V$.
Then as before we have the functor of indecomposables $Q: \HCAlg_{\Z_p}^V \to \Mod_{\Z_p}^{\N,V}$.
\begin{proposition}[Goerss] \label{prop:Hopf-with-V}
The functor $Q$ is cocontinuous, symmetric monoidal and restricts to an equivalence of categories between (on either side) the connected, torsion-free objects.
\end{proposition}
\begin{proof}
Symmetric monoidality is proved as before.
For the equivalence, see \cite[Theorem 2.7]{goerss-dieudonne}.
\end{proof}

\begin{example} \label{ex:M-versch}
For $n \in \N$ we have the object $M_n \in \Mod_{\Z_p}^{\N,V}$ which is given by $\Z_p$ in degrees of the form $n/p^*$ and by $0$ else, with $V$ the identity whenever this makes sense.
Applying the inverse of $Q$, we obtain important Hopf algebras with lift of Verschiebung, still denoted $M_n$.
\end{example}

\begin{example} \label{ex:Mp-hotimes-Mp}
Computations of tensor products of Hopf algebras with lift of Verschiebung are straightforward using the above equivalence.
For example, let us write $M_{p^\infty} := \colim_i M_{p^i}$.
Then \[ M_{p^\infty} \hotimes M_{p^\infty} \wequi \bigoplus_{i > 0} \phi_{p^i + 1}^* M_{p^\infty} \oplus \phi_2^* M_{p^\infty} \oplus \bigoplus_{j > 0} \phi_{p^j + 1}^* M_{p^\infty}. \]
Similarly $M_n \hotimes M_m$ is easily seen to be given by a finite sum of copies of $M_i$, for various $i$.
Symmetric and exterior powers are treated similarly.
\end{example}

\begin{remark}
The similarities between Examples \ref{ex:Mp-hotimes-Mp} and \ref{ex:Ga-hotimes-Ga} are not just superficial.
It is in fact the case that $(M_{p^\infty} \otimes \F_p)/F \wequi \Ga$; see Lemma \ref{lemm:W-mod-F}.
\end{remark}

% COPIED with limited algebraic substitutions from RW-comparison-expanded.tex
% Original portions: wilson.tex lines 1639--1756, 1871--2017, 2057--2146.
\subsection{Homomorphic curves}\label{subsec:RW-moduli}
Let $\scr C$ be an additive, symmetric monoidal $1$-category (in which the tensor product is additive).
By a Hopf ring in $\scr C$, we mean a commutative graded ring object in the cartesian symmetric monoidal category $\cCAlg(\scr C)$.
Equivalently, giving $R \in \cCAlg(\scr C)$ the structure of a Hopf ring is to supply a factorization of \[ \Hom_{\cCAlg(\scr C)}(\ph, R): \cCAlg(\scr C)^\op \to \Set \] through commutative graded rings.
Similarly, given an ordinary commutative graded ring $A$, we can define a Hopf ring under $A$ in $\scr C$ by requiring a factorization of the represented functor through commutative graded $A$-algebras.


\begin{example}
By definition, if $R$ is a Hopf ring under $A$ then the set of (graded) coalgebra maps $\1 \to R$ is an $A$-algebra, and thus contains elements $[a]$ for any $a \in A$.
In the situation of the above remark, elements $a \in A$ (i.e., maps $a: * \to A$) induce the elements $[a] := \1_\scr{C}[a] \in \1_\scr{C}[A]$ (i.e., maps $[a]: \1_\scr{C} \to \1_\scr{C}[A]$) via the functor $\1_\scr{C}[\ph]$.
\end{example}

\begin{example} \label{ex:hopfring-frob}
Consider the case $A=\F_p$, i.e., Hopf rings with $[p]=0$.
Any functor $\cCAlg(\scr C)^\op \to \CAlg_{\F_p/}^\heart$ has a canonical endomorphism, given sectionwise by the Frobenius.
This induces a canonical endotransformation of the identity functor on Hopf rings under $\F_p$.
In more down to earth terms, if $R$ is a Hopf ring with $[p]=0$, then the coalgebra map $R \xrightarrow{\Delta_p} R^{\otimes p} \xrightarrow{m_p} R$ is a Hopf algebra map, where $m$ means the Hopf ring multiplication.
\end{example}

\begin{definition} \label{def:ModLfr}
We write $L$ for the $p$-local Lazard ring, viewed as a commutative graded ring, and $\Mod_L^\fr$ for the symmetric monoidal $1$-category of free graded $L$-modules on generators in non-negative degrees.
For a multiplicative subset $S \subset \Z$, define $\Mod_{S^{-1}L}^\fr$ similarly.
\end{definition}
\begin{remark}\label{rmk:lazard-ring}
We shall use frequently and without further comment that $L$ (as defined above) is also the universal ring carrying a (commutative, 1-dimensional) formal group law \cite{quillen-MU}.
In other words there exist elements $a_{i,j} \in L_{1+i+j}$ such that the formal power series $\sum_{i,j} a_{i,j} x^i y^j \in L\fpsr{x,y}$ is a formal group law with a certain universal property.
\end{remark}

% COPIED from wilson.tex, lines 1670--1678; p-localization is implicit in the coefficient ring L.
\begin{definition} \label{def:A1}
By standard computations in homotopy theory, the $p$-local complex cobordism homology
\[\A^1=\MU_{2*}\C\P^\infty\otimes_{\MU_{2*}}L\]
defines an object of $\cCAlg(\Mod_L^\fr)$.
For $r\ge 0$, write
\[A_r=\A^1_{\le r}=\MU_{2*}\C\P^r\otimes_{\MU_{2*}}L\subset\A^1.\]
\end{definition}
\begin{remark}
The multiplication map $\mu: \C\P^\infty \times \C\P^\infty \to \C\P^\infty$ and the pointing $* \to \C\P^\infty$ in fact propagate $\A^1$ to a Hopf algebra in $\Mod_L^\fr$.
We will however usually only think of $\A^1$ as a coalgebra, and keep track of the multiplication and pointing explicitly.
\end{remark}
\begin{remark}
We can define the above objects without reference to homotopy theory.
The coalgebra $\A^1$ is free over $L$, with basis $\{b_i\}_{i \ge 0}$ where $|b_i| = i$, $\Delta(b_i) = \sum_{i = k+j} b_k \otimes b_j$, and $b_i b_j = {i+j \choose j} b_{i+j}$.
Then the algebra dual to $\A^1$ is the formal power series ring $L\fpsr{c}$ with $b_i$ dual to $c^i$.
The map $\mu$ is most easily described on duals, where $\mu^*: L\fpsr{c} \to L\fpsr{x,y}$ is induced by sending $c$ to the universal formal group law $\sum_{ij} a_{ij} x^i y^j$ (see Remark \ref{rmk:lazard-ring}).
\end{remark}

In what follows, we will routinely find ourselves in situations where certain seemingly infinite sums can (and must) be formed.
The basis of this is the following.
\begin{definition} \label{def:conil}
We call a pointed coalgebra $C \in \cCAlg(\scr C)_{\1/}$ \emph{conilpotent} if there is $N \gg 0$ such that the iterated reduced diagonal \[ C \xrightarrow{\Delta^N} C^{\otimes N} \to (C/\1)^{\otimes N} \] is null.
We call a pointed coalgebra $C \in \cCAlg(\scr C)_{\1/}$ \emph{ind-conilpotent} if we can present $C = \colim_i C_i$ as a filtered colimit in pointed coalgebras, where each $C_i$ is conilpotent.
\end{definition}
% COPIED from wilson.tex, lines 1692--1695, using A_r from Definition def:A1.
\begin{example}
We have $\A^1 \wequi \colim_r A_r$, exhibiting $\A^1$ as ind-conilpotent.
(Indeed $\MU^{2*}\C\P^r\otimes_{\MU_{2*}}L\wequi L[c]/c^{r+1}$, from which conilpotence of $A_r$ easily follows.)
\end{example}

\begin{remark} \label{rmk:pointed-vanishing}
Let $R \in \CAlg(\HCAlg(\scr C))$ and $C = \colim_i C_i$ an ind-conilpotent coalgebra, with each $C_i$ conilpotent.
For any pointed map $x: C_i \to R$ ($R$ pointed by $[0]$), $x^n$ factors through the reduced $n$-fold diagonal (because the multiplication in $R$ induces $R/\1 \otimes R/1 \to R/1$, the pointing being by $[0]$); thus $x^n = [0]$ for $n \gg 0$.
Since also $\Hom(C, R) \wequi \lim_i \Hom(C_i, R)$, for $x \in \Hom(C, R)$ we can make sense of infinite sums like $\sum_i x^i$ (since they are finite when restricted to any $C_i$).
\end{remark}

We will now work towards defining the Ravenel--Wilson Hopf rings, via the moduli problem they corepresent.
This is done in the context of the following data:
\begin{itemize}
\item A commutative graded formal group law $L \to A$.
\item A strongly additive symmetric monoidal functor $F: \Mod_L^\fr \to \scr C$ (where $\scr C$ is some strongly additive symmetric monoidal $1$-category).
\end{itemize}
\begin{remark} \label{rmk:FA1-indconil}
Note that $F(\A^1) \wequi \colim_i F(A_i)$, since the colimit on underlying objects (ignoring the coalgebra structures) is an infinite sum (and thus preserved by $F$, by assumption).
In other words, $F(\A^1)$ is ind-conilpotent.
\end{remark}
\begin{definition} \label{def:homomorphic-curve}
In the above situation, for any Hopf ring $R$ in $\scr C$ under $A$ we define the \emph{set of homomorphic curves in $R$} as the set of all coalgebra maps $\gamma: F(\A^1) \to R^1$ such that:
\begin{itemize}
\item $\gamma$ is pointed.
\item Writing $\mu: \A^1 \times \A^1 \to \A^1$ for the multiplication map and $p_i: \A^1 \times \A^1 \to \A^1$ for the projections, we have \[ \mu^*(\gamma) = \sum_{i,j} [a_{ij}] p_1^*(\gamma)^i p_2^*(\gamma)^j. \]
  Here on the right hand side, the multiplication uses the multiplication in $R$, the sum uses the addition in $R$, and the sum is defined by Remarks \ref{rmk:pointed-vanishing} and \ref{rmk:FA1-indconil}.
\end{itemize}

If this functor is corepresentable by a Hopf ring in $\scr C$ under $A$ (e.g., $\scr C$ presentable), we denote by $\RW(\scr C, A)$ the corepresenting object.
\end{definition}
\begin{remark} \label{rmk:construct-RW}
In other words, there exists at most one Hopf ring $\RW(\scr C, A)$ under $A$ in $\scr C$ (up to unique isomorphism) such that $A$-algebra homomorphisms from $\RW(\scr C, A)$ to $R$ are the same as homomorphic curves in $R$.
In fact, if $\scr C$ has enough colimits, we can build a pushout of the form
\begin{equation*}
\begin{CD}
F(\A^2) @>{\alpha}>> F(\A^1) \\
@VVV           @VVV  \\
\1_\scr{C}[L] @>>> \RW(\scr C, A).
\end{CD}
\end{equation*}
Here $F(\ph)$ denotes the free Hopf ring on a pointed coalgebra placed in degree $1$, and the map $\alpha$ corepresents the map \[ \gamma \mapsto \mu^*(\gamma) - \sum_{i,j} [a_{ij}] p_1^*(\gamma)^i p_2^*(\gamma)^j. \]
Then the bottom right corner corepresents homomorphic curves, essentially by construction.
\end{remark}


\begin{remark}[Additive curves] \label{rmk:additive-curves}
Let $F: \Mod_L^\fr \to \scr C$ be an additive symmetric monoidal functor.
Then the Hopf algebra $(\A^1, \mu)$ is sent by $F$ to another Hopf algebra $(F\A^1, F\mu)$.
It is possible that $(F\A^1, F\mu)$ is isomorphic to $F\Ga$.
(E.g., take $L \to \F_p$ and the induced functor $F: \Mod_L^\fr \to \Vect_{\F_p}$.)
In this case homomorphic curves in $\scr C$ are usually called \emph{additive curves}.
\end{remark}


% COPIED from original wilson.tex lines 1871--2146; source comments preserved.
\subsection{Curves in formal groups}\label{subsec:curves}
In preparation for the next section, here we discuss some further structures in $\Hopf_{\F_p}$.
We recalled above the functors $\phi_n^*$ and the endotransformations $F,V$.

\begin{definition}
By a \emph{formal group} we shall mean a functor \[ X: \cCAlg(\Vect_{\F_p}^\N)^{\mathrm{fd},\cn,\op}_{\F_p/} \to \Ab \] from (the opposite of) pointed, connected, finite dimensional (over $\F_p$) coalgebras to abelian groups, such that there exists a pointed coalgebra $\tilde X$ and an isomorphism \[ X|_{\Fun(\cCAlg(\Vect_{\F_p}^\N)^{\mathrm{fd},\cn, \op}_{\F_p/}, \Set)} \wequi \Hom(\ph, \tilde X). \]
\end{definition}

One class of examples arises as follows.
\begin{example} \label{ex:fg-1}
Let $H \in \Hopf_{\F_p}$.
Define $H^\comp(C) = \Hom_{\cCAlg(\Vect_{\F_p}^\N)_{\F_p/}}(C, H)$.
This inherits an abelian group structure from $H$, and so is a formal group.
\end{example}

Now let $X$ be a formal group.
We call \[ X_n = X(\phi_n^* \A^1) \] the set of \emph{pointed curves of degree $n$ in $X$}.

\subsubsection{Formal groups and Hopf algebras}
Our definition of formal groups has mostly psychological utility, as the following result shows.
\begin{lemma} \label{lemm:HCAlg-vs-FG}
The functor \[ (\ph)^\comp: \Hopf_{\F_p}^\cn \to (\text{(formal groups)}) \] is an equivalence of categories.
\end{lemma}
\begin{proof}
Let us call a functor
\[X: \cCAlg(\Vect_{\F_p}^\N)^{\mathrm{fd},\cn,\op}_{\F_p/} \to \Set\]
representable by a pointed coalgebra $\tilde X$ a \emph{pointed formal scheme}.
We have
\[(\text{formal groups}) = \Ab(\text{pointed formal schemes}).\]
By the fundamental theorem of coalgebras \cite[Corollary 2.2.2]{sweedler-hopf} we have
\[\cCAlg(\Vect_{\F_p}^\N)^\cn_{\F_p/}
 \wequi \Ind(\cCAlg(\Vect_{\F_p}^\N)^{\mathrm{fd},\cn}_{\F_p/}).\]
Consequently the restricted Yoneda functor
\[\cCAlg(\Vect_{\F_p}^\N)^\cn \to (\text{formal schemes})\]
is fully faithful.
It will suffice to prove that it is essentially surjective.
I.e., given a pointed coalgebra $\tilde X$, we need to find a pointed connected coalgebra $X$ such that $X^\comp \wequi \tilde X^\comp$.
We can take $X$ to be the image of \[ \bigvee_{C \to \tilde X} C \to \tilde X, \] where the wedge is over all morphisms of pointed coalgebras $C \to \tilde X$ with $C$ finite dimensional.\NB{details?}
\end{proof}

Denote by $\Free$ the left adjoint of the forgetful functor $\Hopf_{\F_p} \to \cCAlg(\Vect_{\F_p}^\N)_{\F_p/}$.
The Hopf algebra $\Free\A^1$ is called the \emph{big Witt Hopf algebra}.
As a ring we have \[ \Free\A^1 = \F_p[a_1, a_2, \dots]; \] in particular $\Free\A^1$ is connected.
In fact, $(\Free\A^1)^\comp$ is the free formal group on $\A^1$, in the following sense.

\begin{corollary} \label{cor:Free-FG}
Let $X$ be a formal group and $\gamma \in X_n$.
Then $\gamma$ extends uniquely to a morphism of formal groups $(\Free\phi_n^*\A^1)^\comp \to X$.
\end{corollary}
\begin{proof}
Immediate from Lemma \ref{lemm:HCAlg-vs-FG} and connectivity of $\Free\phi_n^*\A^1$.
\end{proof}

\begin{corollary} \label{cor:fg-p-local}
All formal groups are $p$-local, i.e., for any $m$ with $(m,p)=1$, the multiplication by $m$ map is an isomorphism on any formal group.
\end{corollary}
\begin{proof}
By Lemma~\ref{lemm:HCAlg-vs-FG}, it suffices to consider a connected graded Hopf algebra. In any fixed positive degree, a sufficiently high iterate of $V$ vanishes for degree reasons. Thus $[p]=FV$ acts locally nilpotently by Lemma~\ref{lemm:FV-basics}, and hence multiplication by any integer prime to $p$ is invertible.
\end{proof}

\subsubsection{Frobenius and Verschiebung for curves}
We will now construct the Frobenius and Verschiebung operators on the sets $X_n$.

\begin{definition}
Let $X$ be a formal group and $\gamma: \phi_n^* \A^1 \to X$ a pointed curve of degree $n$.
Form \[ \phi_n^* \A^m \xrightarrow{p_1^*\gamma + \dots + p_m^*\gamma} X, \] which is symmetric in the $m$ coordinates, and hence factors uniquely through the quotient coalgebra $\phi_n^*\A^m/\Sigma_m$.\footnote{The existence of $\tilde X$ implies both that the factorization exists and is unique.}
Define the \emph{Frobenius operator} $\frob_m(\gamma) \in X_{mn}$ as the composite \[ \phi_{mn}^*\A^1 \to \phi_n^*\A^1 \htimes \dots \htimes \phi_{nm}^* \A^1 \wequi \phi_n^*\A^m/\Sigma_m \xrightarrow{p_1^*\gamma + \dots + p_m^*\gamma} X, \] where the equivalence is induced by the symmetric function theorem.\footnote{In other words, passing to duals, an isomorphism of coalgebras $\phi_n^*\A^1 \htimes \dots \htimes \phi_{nm}^* \A^1 \wequi \phi_n^*\A^m/\Sigma_m$ is the same as an isomorphism of rings $\F_p[X_n, \dots, X_{mn}] \wequi \F_p[Y_n^{(1)}, \dots, Y_n^{(m)}]^{\Sigma_m}$, which is supplied by the classical symmetric function theorem.}

We also define the \emph{Verschiebung operators} given by \[ \ver_m: X_{mn} \to X_{n}, (\phi_{mn}^*\A^1 \xrightarrow{\gamma} X) \mapsto (\phi_n^* \A^1 \to \phi_{mn}^* \A^1 \xrightarrow{\gamma} X). \]
Here the map $\phi_n^* \A^1 \to \phi_{mn}^* \A^1$ is dual to $\F_p[t_{nm}] \to \F_p[t_n], t_{mn} \mapsto t_n^m$.
\end{definition}

\begin{remark}
Essentially by definition, $\phi_n^* \Free\A^1$ corepresents the functor \[ \Hopf_{\F_p} \to \Ab, \quad H \mapsto H^\comp_n. \]
The operations $\frob_m$, $\ver_m$ on pointed curves thus correspond to operations \[ F_m: \phi_{mn}^* \Free\A^1 \to \phi_n^*\Free\A^1 \text{ and } V_m: \phi_n^*\Free\A^1 \to \phi_{mn}^* \Free\A^1. \]
\end{remark}

% COPIED from wilson.tex, lines 2018--2028; source comments retained.
\begin{lemma} \label{lemm:FV-on-W}
On $\Free\A^1$ have $F_p = F$ and $V_p = V$.
\end{lemma}
\begin{proof}
The map $F_p$ is obtained as the composite \[ \Free \phi_p^* \A^1 \to \Free(\A^p/\Sigma_p) \to (\Free \A^1)^{\htimes p}/\Sigma_p \to \Free \A^1, \] where the first map comes from the symmetric function theorem, the second map is a categorical comparison map, and the third is the multiplication.
Unwinding the definitions, this is precisely $F$.\NB{details?}
Since $V$ is already defined for coalgebras and functorial, in order to prove that $V=V_p$ for $\Free\A^1$, it suffices to prove the same for $\A^1$ itself.
In other words, we must show that $V$ on $\A^1$ is dual to the Frobenius.
This is true essentially by definition.
\end{proof}

% COPIED from wilson.tex, lines 2030--2034; source comments retained.
\begin{example} \label{ex:lift-FreeA1}
It follows that $\Free\A^1$ lifts to a torsion-free Hopf algebra over $\Z_p$, with a lift of Verschiebung given by $V_p$ (constructed the same way over $\Z_p$ as over $\F_p$).
We can analyze this lifted Hopf algebra as in \S\ref{subsec:lift-of-v}, i.e., in terms of the action of $V$ on its indecomposables.
Using \cite[(13.4)]{hazewinkel-witt-vectors}, we see that the indecomposables are given, in the notation of Example \ref{ex:M-versch}, by $M_\infty := \colim_n M_n$.
\end{example}

\begin{remark} \label{rmk:symm-hazewinkel}
The Hopf algebra $\Free\A^1$ is denoted by $\mathrm{Symm}$ in \cite{hazewinkel-witt-vectors}.
(See in particular (9.4) and (9.5) in \emph{loc. cit.})
Moreover, unwinding the definitions, the maps $F_m$ and $V_m$ we have defined coincide with Hazewinkel's \cite[\S13]{hazewinkel-witt-vectors}.
This is useful because \emph{loc. cit.} establishes numerous properties of these operations.
(Hazewinkel does not keep track of gradings the way we do, but most desired properties can be deduced nonetheless.)
\end{remark}

These operators satisfy the following relations.
\begin{lemma} \label{lemm:frob-ver-identities}
The operations $\frob_m$ and $\ver_m$ are natural in $X$ and satisfy
\begin{align*}
\frob_1 &= \ver_1 = 1 \\
\frob_m\ver_m &= m \\
\frob_m\frob_m' &= \frob_{m m'} \\
\ver_m\ver_{m'} &= \ver_{m m'} \\
\frob_m\ver_{m'} &= \ver_{m'}\frob_m \quad (m',m)=1 \\
\ver_p\frob_p &= p,
\end{align*}
whenever they are defined.
\end{lemma}
Note that $\ver_m$ can only be applied to curves of degree divisible by $m$.
So the second equation above claims that the endomorphism $\frob_m \ver_m$ of $X_{nm}$ is given by multiplication by $m$ for every $X, n$.
If $m$ does not divide $n$, then $\ver_m$ on $X_n$ is not defined, not even as zero, and we are not claiming that multiplication by $m$ on $X_n$ is zero.
\begin{proof}
By Corollary~\ref{cor:Free-FG}, this reduces to the universal case $X=(\Free\A^1)^\comp$, i.e., the big Witt vectors.
The relations in this context are proved in \cite[Theorem 13.48]{hazewinkel-witt-vectors}, except for the last one, which is a consequence of Lemma~\ref{lemm:FV-on-W}.
\end{proof}

\begin{definition}
We call a pointed curve $\gamma \in X_n$ \emph{$p$-typical} if $\frob_\ell \gamma = 0$ for all primes $\ell \ne p$.
\end{definition}

With this at hand, we can now construct the \emph{$p$-typicalization operator}.
\begin{lemma} \label{lemm:p-typicalization}
Let $\mu$ be the Möbius function.\footnote{Thus $\mu(n) = 0$ if $n$ is divisible by a square $>1$ and $\mu(n) = (-1)^k$ if $n$ is a product of $k$ distinct primes.}
Suppose that $X$ is a formal group.
The operator
\[
\wittprojection_p = \sum_{(m,p)=1} \frac{\mu(m)}{m} \ver_m \frob_m
\]
is well-defined and has the following properties:
\begin{enumerate}
\item For all $\gamma$, the curve $\wittprojection_p \gamma$ is $p$-typical;
\item If $\gamma$ is $p$-typical then $\wittprojection_p \gamma=\gamma$;
\item $\wittprojection_p \wittprojection_p =\wittprojection_p$.
\end{enumerate}
\end{lemma}
\begin{proof}
We first check that the operator is well-defined.
The division by $m$ expressions make sense by Corollary \ref{cor:fg-p-local}.
Let $A_r \subset \A^1$ be the subcoalgebra spanned by $b_0,\dots,b_r$, so that $\A^1 = \colim_r A_r$.
We must show that when restricting along $\phi_m^* A_r \to \phi_m^* \A^1$, the sum defining the operator becomes finite.
For this, it will be enough to show that given $r$ there exists $N$ such that for all $n>N$ and any pointed curve $\gamma: \phi_{mn}^* \A^1 \to X$, the restriction of $\ver_n \gamma$ to $\phi_m^* A_r$ is zero.
Equivalently, using the definition of $\ver_n$, we want that the composite $\phi_m^* A_r \to \phi_m^* \A^1 \to \phi_{mn}^* \A^1$ (where the last map is the $n$-th power) is null.
This is dual to the composite $\F_p[X_{mn}] \xrightarrow{X_{mn} \mapsto X_m^n} \F_p[X_m]/X_m^{r+1}$, so $N=r$ works.


(1) Fix a prime $\ell \ne p$; we shall show that $\frob_\ell \wittprojection_p = 0$.
Let $S = \{m \mid (m,p) = 1 = (m,\ell)\}$.
Then \[ \wittprojection_p = \sum_{m \in S} \frac{\mu(m)}{m} \ver_m \frob_m + \sum_{m \in S} \frac{\mu(\ell m)}{\ell m} \ver_{\ell m} \frob_{\ell m}. \]
Using that $\mu(\ell m) = -\mu(m)$ for $m \in S$ and the identities in Lemma \ref{lemm:frob-ver-identities}, the claim follows quickly.

(2) Immediate.

(3) Clear from (1) and (2).
\end{proof}

% The Witt Hopf algebra is defined after the p-typicalization operator.
\subsubsection{Witt Hopf algebra} \label{app:rw-witt-hopf}
Let $H\in\Hopf_{\F_p}$. Applying the $p$-typicalization operator of Lemma~\ref{lemm:p-typicalization} to $H^\comp_1$ defines a retract of the functor corepresented by $\Free\A^1$.
\begin{definition}\label{def:W}
The corresponding summand of $\Free\A^1$ is denoted by $W$ and called the \emph{Witt Hopf algebra}. As an algebra it is polynomial, with generators $u_i$ in degrees $p^i$ (after normalizing the first degree to one).
\end{definition}

% COPIED from wilson.tex, lines 2042--2055; source comments retained.
\begin{remark} \label{rmk:W-lift}
Comparing \cite[Theorem 14.21]{hazewinkel-witt-vectors} with our construction of $p$-typicalization, we see that our definition of $W$ agrees with the usual (Hopf algebra representing) $p$-adic Witt vectors.
As in Example \ref{ex:lift-FreeA1}, $W$ lifts to $\HCAlg_{\Z_p}^V$ and, being torsion-free, the lift is determined by its indecomposables with their $V$-action.
In the notation of Example \ref{ex:Mp-hotimes-Mp}, it is given by $M_{p^\infty}$ (use \cite[\S14.8, (13.4)]{hazewinkel-witt-vectors}).
\end{remark}

\begin{lemma} \label{lemm:W-mod-F}
$F: \phi_p^* W \to W$ is injective and $W/F \wequi \Ga$, via the canonical map.
\end{lemma}
\begin{proof}
That $F$ is injective on $W$ follows from the same statement for $\Free\A^1$, which is clear because $\Free\A^1$ is a polynomial ring.
Since $F=0$ on $W/F$ as well as $\Ga$, by Proposition \ref{prop:indecomp-charp} it suffices to check that $W/F \to \Ga$ induces an isomorphism on indecomposables.
This is true by Remark \ref{rmk:W-lift} and Example \ref{ex:Ga-hotimes-Ga}.
\end{proof}

\begin{remark}\label{app:rw-projective}
The Hopf algebras $M_n\otimes \F_p$ above are projective generators of the additive category of connected graded Hopf algebras \cite{goerss-dieudonne,MR281709}. Hence $W$ is a filtered colimit of projectives, and its symmetric powers are direct sums of twists of $W$ (compare Example~\ref{ex:Mp-hotimes-Mp}). In particular, the symmetric powers of $W$ have flat underlying objects for $\hotimes$.
\end{remark}
Let us also observe the following.
\begin{lemma} \label{lemm:additive-curves-frob}
A pointed curve $\gamma \in X_n$ corresponds to a morphism of formal groups $\phi_n^* \Ga^\comp \to X$ if and only if $\frob_q \gamma = 0$ for all primes $q$.
\end{lemma}
\begin{proof}
In light of Lemma \ref{lemm:HCAlg-vs-FG} (identifying formal groups and connected Hopf algebras), an equivalent formulation is that the sequence \[ \bigotimes_q \Free \phi_q^* \A^1 \xrightarrow{F_{(\ph)}} \Free\A^1 \to \Ga \] is right exact.
For the $p$-typical part, this is Lemma \ref{lemm:W-mod-F} above.
\end{proof}


\subsection{An alternative presentation of the Ravenel--Wilson Hopf ring}\label{subsec:RW-p}
In this section we consider the case $\scr C = \Vect_{\F_p}^\N$, $A = L$.
We put $\RW_{(p)} := \RW(\Vect_{\F_p}^\N, L)$, and we shall construct an alternative presentation for this ring.

Here is a new example of a formal group.
\begin{example}\label{ex:fg-2}
Let $R \in \CAlg(\Hopf_{\F_p}^\N)_{\F_p[L]/}$.
Then the set $(R^1)^\comp(C)$ defined in Example \ref{ex:fg-1} acquires a new abelian group structure denoted $+^F$, where \[ \gamma_1 +^F \gamma_2 = \sum_{ij} [a_{ij}] \gamma_1^i \gamma_2^j. \]
This sum is well-defined by conilpotence (see Remark \ref{rmk:pointed-vanishing}; note that a connected finite-dimensional coalgebra is always conilpotent).
\end{example}
Via Examples \ref{ex:fg-1} and \ref{ex:fg-2}, we have thus found \emph{two} formal group structures on the set-valued functor $(R^1)^\comp$.
To distinguish them, we denote the operations by $+$ (Example \ref{ex:fg-1}) and $+^F$ (Example \ref{ex:fg-2}).
Correspondingly we obtain two notions of Frobenius operators, which we write $\frob_q$ and $\frob_q^F$, and we speak of $p$-typical and $p^F$-typical curves.

\begin{theorem} \label{thm:alternative-presentation}
\begin{enumerate}
\item There is a canonical isomorphism between the two functors \[ \CAlg(\Hopf_{\F_p}^\Z)_{L/} \to \Set \] sending $R$ to the $p$-typical curves, respectively $p^F$-typical curves (in $(R^1)^\comp$).
\item There is a pushout square in $\CAlg(\Hopf_{\F_p}^\Z)_{L/}$ of the form
\begin{equation*}
\begin{CD}
\Sym_L \phi_p^* W @>>> L_{(p)} \\
  @V\frob_p^FVV        @VVV \\
\Sym_L W @>>> \RW_{(p)}.
\end{CD}
\end{equation*}
  Here $W$ is placed in degree $1$ and $\Sym_L W$ corepresents, via (1), the functor of $p^F$-typical curves.
  Under this isomorphism, the map $\frob_p^F$ corepresents the operator of the same name.
\end{enumerate}
\end{theorem}
Let us elaborate slightly on part (2).
The functor $\Hopf_{\F_p} \to \Ab$, $H \mapsto \wittprojection_p H^\comp_1$ is corepresented by $W$, essentially by definition.
Consequently the Hopf ring $\Sym_L W$ corepresents the functor sending $R$ to the $p$-typical curves in $R^1$.
By part (1), this functor is isomorphic (but not equal) to the functor sending $R$ to the $p^F$-typical curves in $R^1$.
The operation $\frob_p^F$ on $p^F$-typical curves thus corresponds to an endomorphism (up to twisting) of the functor corepresented by $\Sym_L W$, whence an endomorphism (up to twisting) of this object.
This endomorphism is what appears in the left hand vertical map of the diagram.
\begin{proof}
Assuming (1), (2) is immediate from Lemma \ref{lemm:additive-curves-frob}.
We thus prove (1).
On $X=(R^1)^\comp$, we have two formal group structures denoted $+$ and $+^F$.
We consequently obtain two notions of $p$-typicalization, denoted $\wittprojection_p$ and $\wittprojection_p^F$.
We claim that more precisely, $\wittprojection_p$ induces a bijection from $p^F$-typical curves to $p$-typical curves.
Write $I \subset L_{(p)}$ for the augmentation ideal.
The two representing Hopf rings are positive, as follows from their presentations by free Hopf rings on the curve coalgebra. Since the generators of $I$ have strictly negative external degrees, in each fixed bidegree sufficiently high powers of $I$ vanish on connected components. Thus the two representing Hopf rings are $I$-adically complete.\gpt{Replaces references to abstract positivity and completion results.}
It thus suffices to prove the claim for $I$-adically complete Hopf rings (under $L_{(p)}$).
Let $R$ be one such.

\emph{Injectivity}:
Let $\gamma_1, \gamma_2: \A^1 \to R^1$ be $p^F$-typical curves with $\wittprojection_p \gamma_1 = \wittprojection_p \gamma_2$.
We need to show that $\gamma_1 = \gamma_2$, and by completeness, it will suffice to show that $\gamma_1$ and $\gamma_2$ have the same image in $R/I^n R$ for every $n$.
We do this by induction, the case $n=0$ being clear.
%The $L_{(p)}$-module map $I^n \to L_{(p)}$ induces $\F_p[I^n] \to \F_p[L]$ and hence \[ \F_p[I^n] \hotimes_{\F_p[L]} R \to \F_p[L] \hotimes_{\F_p[L]} R \wequi R \in \Mod_{\F_p[L]}(\Hopf_{\F_p}^\Z); \] denote its image by $I^n R$.
We have an exact sequence \[ \F_p \to I^nR/I^{n+1}R \to R/I^{n+1}R \to R/I^nR \to \F_p \in \Mod_{\F_p[L]}(\Hopf_{\F_p}^\Z). \]
Applying the (trivially!) left exact functor $\Hom(\A^1, (\ph)^1)$ we get \[ \begin{aligned} 0 &\to \Hom(\A^1,(I^nR/I^{n+1}R)^1) \\ &\to \Hom(\A^1,(R/I^{n+1}R)^1) \to \Hom(\A^1,(R/I^n R)^1). \end{aligned} \]
By induction, the images of $\gamma_1, \gamma_2$ in the right hand group agree, and consequently we find $\delta \in \Hom(\A^1, (I^nR/I^{n+1}R)^1)$ with \[ \overline{\gamma_1} = \overline{\gamma_2} + \delta \in \Hom(\A^1, (R/I^{n+1}R)^1). \]
Note that $R/I^{n+1}R$ is still a ring, and so curves in $R/I^{n+1}$ can still be added using the formal group law (producing an alternative group structure on $\Hom(\A^1, (R/I^{n+1}R)^1)$ denoted $+^F$).
We then have $\overline{\gamma_2} +^F \delta = \overline{\gamma_2} + \delta$ since $[a_{ij}]\delta = 0$ for any $a_{ij} \in I$.
Consequently $\delta = \overline{\gamma_1} -^F \overline{\gamma_2}$ is $p^F$-typical (by construction, $\wittprojection_p^F$ is a homomorphism for $+^F$).
On the other hand $\wittprojection_p(\delta) = \wittprojection_p(\overline{\gamma_1} - \overline{\gamma_2}) = 0$, since $\wittprojection_p$ is a homomorphism for the usual $+$ and $\gamma_1, \gamma_2$ are supposed to have the same $p$-typicalization.
We claim that $\wittprojection_p(\delta) = \wittprojection_p^F(\delta)$.
Assuming this, since $\delta$ was shown to be $p^F$-typical, we deduce that $\delta=0$ as needed.
To prove the claim, recall that the set-valued functor $((R/I^{n+1})^1)^\comp$ carries two formal group structures, one coming from the usual Hopf algebra structure and one involving the $[a_{ij}]$.
This functor has a subset, namely $((I^nR/I^{n+1}R)^1)^\comp$.
This is a subgroup for either formal group structure (namely the kernel of $R/I^{n+1}R \to R/I^nR$), and the two formal group structures on the subfunctor in fact coincide (because, as above, given $x: C \to I^nR/I^{n+1}R$ we have $[a_{ij}]x = 0$ for any $a_{ij} \in I$).
This proves that $\wittprojection_p = \wittprojection_p^F$ on the subfunctor, as needed.

\emph{Surjectivity}:
Let $\gamma: \A^1 \to R^1$ be $p$-typical.
To simplify notation, given $a, b: \A^1 \to R^1$, we shall say that $a=b \pmod{I^n}$ if $a$ and $b$ have the same image in $\Hom(\A^1, (R/I^nR)^1)$.
We shall construct a sequence of $p^F$-typical curves $\tilde \gamma_n$ such that for each $n$, (1) $\tilde\gamma_n = \tilde \gamma_{n-1} \pmod{I^{n-1}}$  and (2) $\wittprojection_p \tilde\gamma_n = \gamma \pmod{I^n}$.
Assuming this, the $\tilde \gamma_n$ define an element \[ \tilde \gamma \in \Hom(\A^1, R^1) \wequi \lim_n \Hom(\A^1, (R/I^nR)^1). \]
By construction, $\wittprojection_p^F(\tilde \gamma) = \tilde \gamma \pmod{I^n}$ for every $n$, i.e. $\wittprojection_p^F(\tilde \gamma)$ and $\tilde \gamma$ have the same image in $\Hom(\A^1, (R/I^nR)^1)$ for every $n$, so agree; in other words, $\tilde\gamma$ is $p^F$-typical.
Also by construction, $\wittprojection_p(\tilde \gamma) = \gamma \pmod{I^n}$ for every $n$, so $\wittprojection_p(\tilde \gamma) = \gamma$.
It thus suffices to construct inductively the $\tilde\gamma_n$; we set $\tilde\gamma_0 = [0]$.
We put \[ \begin{aligned} \tilde\gamma_{n+1} &= \wittprojection_p^F\bigl(\tilde \gamma_n -^F \\ &\qquad \wittprojection_p^F \wittprojection_p (\tilde \gamma_n - \gamma)\bigr). \end{aligned} \] and verify (1) and (2).
By induction, $\wittprojection_p(\tilde \gamma_n - \gamma) = \wittprojection_p(\tilde \gamma_n) - \wittprojection_p(\gamma) = 0 \pmod{I^n}$.
This implies that $\tilde\gamma_{n+1} = \wittprojection_p^F(\tilde \gamma_n -^F [0]) = \tilde \gamma_n \pmod{I^n}$, which is (1).
As in the proof of injectivity, this also implies that \[ (*) \quad \wittprojection_p^F \wittprojection_p(\tilde \gamma_n - \gamma) = \wittprojection_p \wittprojection_p(\tilde \gamma_n - \gamma) \pmod{I^{n+1}} \] and \[ (**) \quad a -^F \wittprojection_p(\tilde \gamma_{n}-\gamma) = a - \wittprojection_p(\tilde \gamma_{n}-\gamma) \pmod{I^{n+1}} \] for any $a: \A^1 \to R^1$.
Now we compute
\begin{align*}
\wittprojection_p\tilde\gamma_{n+1} &= \wittprojection_p\wittprojection_p^F\big( \tilde\gamma_{n} -^F \wittprojection_p^F\wittprojection_p(\tilde \gamma_{n}-\gamma)\big) \pmod{I^{n+1}} \\
&=\wittprojection_p(\wittprojection_p^F \tilde\gamma_{n} -^F \wittprojection_p^F\wittprojection_p^F\wittprojection_p(\tilde \gamma_{n}-\gamma)\big) \pmod{I^{n+1}} \\
&=\wittprojection_p(\tilde\gamma_{n} -^F \wittprojection_p^F\wittprojection_p(\tilde \gamma_{n}-\gamma)\big) \pmod{I^{n+1}} \\
&\stackrel{(*)}{=}\wittprojection_p(\tilde\gamma_{n} -^F \wittprojection_p\wittprojection_p(\tilde \gamma_{n}-\gamma)\big) \pmod{I^{n+1}} \\
&=\wittprojection_p(\tilde\gamma_{n} -^F \wittprojection_p(\tilde \gamma_{n}-\gamma)\big) \pmod{I^{n+1}} \\
&\stackrel{(**)}{=}\wittprojection_p(\tilde\gamma_{n} - \wittprojection_p(\tilde \gamma_{n}-\gamma)\big) \pmod{I^{n+1}} \\
&=\wittprojection_p\tilde\gamma_{n} - (\wittprojection_p\tilde \gamma_{n}-\gamma) \pmod{I^{n+1}} \\
&=\gamma \pmod{I^{n+1}}.
\end{align*}
\end{proof}



\subsection{The derived Ravenel--Wilson Hopf ring}\label{app:rw-bar}
We recall the definition of the map to be used in the bar construction. For any Hopf ring $T$ under $L$, the pointed curves in $T^1$ have two formal group laws: the one induced by the Hopf-algebra addition, and the one obtained by applying the universal formal group law $\mathcal F(X,Y)=\sum a_{ij}X^iY^j$ using the Hopf-ring product. We denote the associated Frobenius operations by $\frob_p$ and $\frob_p^{\mathcal F}$, respectively. The operation $\frob_p^{\mathcal F}$ on the $p^{\mathcal F}$-typical universal curve determines a morphism
\[
 f=\frob_p^{\mathcal F}:\Sym_L(\phi_p^*W)\longrightarrow \Sym_L(W).
\]
The augmentation $\epsilon:\Sym_L(\phi_p^*W)\to L$ kills the positive external degrees. Put
\[
 A=\Sym_L(\phi_p^*W),\qquad B=\Sym_L(W),\qquad
 R=B\hotimes_A L.
\]
By Theorem~\ref{thm:alternative-presentation}, we have $R\simeq\RW_{(p)}$.

\begin{definition}\label{def:rw-bar}
Define the derived Ravenel--Wilson Hopf ring $\dRW_{(p)}$ to be the simplicial Hopf ring given by the two-sided bar construction
\[
 (\dRW_{(p)})_n=B\hotimes_L A^{\hotimes_L n}\hotimes_L L.
\]
The first and last face maps are induced by $f$ and $\epsilon$; the remaining face maps multiply adjacent $A$-factors. Degeneracies insert the unit. There is a canonical augmentation $\dRW_{(p)}\to R$.
\end{definition}

\begin{theorem}\label{thm:rw-bar-exact}
The augmentation $\dRW_{(p)}\to R$ induces an isomorphism on zeroth homology and $\hH_i(\dRW_{(p)})=0$ for $i>0$.
\end{theorem}

\begin{remark}\label{rmk:rw-bar-derived}
The degree-zero statement follows formally from the bar faces: $\hH_0(\dRW_{(p)})=B\hotimes_A L$.
In the language of derived Hopf rings, the above bar construction computes the derived pushout
\[
\LSym_L^\hotimes(W)\hotimes_{\LSym_L^\hotimes(\phi_p^*W)}L.
\]
\end{remark}

The remainder of this appendix is dedicated to proving Theorem \ref{thm:rw-bar-exact}.

\subsection{The special fiber}\label{app:rw-fiber}

% COPIED from wilson.tex, lines 2177--2265; source comments retained.
\begin{lemma} \label{lemm:filter-derived-RW}
\begin{enumerate}
\item Write $\tilde y: \A^1 \to W \to (\Sym_{L_{(p)}} W)^1 \to \dRW_{(p)}^1$ for the canonical ($[p]$-typical) coordinate.
  Then the composite $\A^1 \xrightarrow{\tilde y} \dRW_{(p)}^1 \to (\dRW_{(p)} \hotimes_L \Z_{(p)})^1$ is an equivalence of coalgebras (so in particular the right hand side is concentrated in homological degree zero) compatibly with addition, i.e., of Hopf algebras.
\item Consequently $(\dRW_{(p)} \hotimes_L \F_p)^1 \wequi \cof(\Ga \xrightarrow{p} \Ga) \wequi \Ga \htimes \Ga[1]$.
  Denote the copy of $\Ga$ in degree zero by $y$ and the one in homological degree $1$ by $x$.
\item We have \[ \hH_* (\dRW_{(p)} \hotimes_L \F_p) \wequi \Sym(\Ga) \hotimes \Lambda(\Ga), \] where the generator $\Ga \wequi \Sym^1(\Ga)$ is $y$ and the generator $\Ga \wequi \Lambda^1(\Ga)$ is $x$.
\end{enumerate}
\end{lemma}
\begin{proof}
By Definition~\ref{def:rw-bar}, the specialization $\dRW_{(p)} \hotimes_L \Z_{(p)}$ is obtained as the bar construction corresponding to the pushout
\begin{equation*}
\begin{CD}
\Sym_{\Z_{(p)}} \phi_p^* W @>>> \Z_{(p)} \\
  @V\frob_p^FVV        @VVV \\
\Sym_{\Z_{(p)}} W @>>> \dRW_{(p)} \hotimes_L \Z_{(p)}.
\end{CD}
\end{equation*}
Because the formal group law on $\Sym_{\Z_{(p)}} W$ is additive, the map $\frob_p^F$ is just $\Sym(\frob_p)$\NB{details?}, i.e., $\Sym(F)$ (Lemma \ref{lemm:FV-on-W}).
Note that $\Sym^0_{\Z_{(p)}}(X) = \F_p[\Z_{(p)}]$ and $\Sym^i_{\Z_{(p)}}(X) = \Sym^i(X)$ whenever $X$ is connected.

(1) Restricting the bar construction defining $\dRW_{(p)} \hotimes_L \Z_{(p)}$ to degree $1$ produces a bar construction modelling $\cof(F: \phi_p^* W \to W)$.
This is $\Ga$ by Lemma \ref{lemm:W-mod-F}.

(2) Immediate.

(3)
We use the filtration by iterates of Frobenius on the normalized simplicial Hopf-algebra complex. Its associated graded has a spectral sequence, converging on connected components by the positive internal grading.\gpt{Termwise Frobenius filtration replaces abstract derived filtration.}
Moreover, this construction is multiplicative.
The degree zero part of $\dRW_{(p)} \hotimes_L \F_p$ is clearly just $\F_p[\F_p]$, so we need only determine the connected part.
We will thus determine $\hH_*(\dRW_{(p)} \hotimes_L \F_p)$ by studying the spectral sequence of the filtration by Frobenius on $\dRW_{(p)} \hotimes_L \F_p$.

\emph{Step 1. Determining the associated graded.}
Since $F$ commutes with $\phi_p^*$, the associated graded in each (filtration) degree is just a twist of the associated graded in degree zero.
It thus suffices to determine associated graded in degree zero, i.e. to determine $\hH_*(\dRW_{(p)} \hotimes_L \F_p/F)$.
We will first do this for $\hH_*(\dRW_{(p)} \hotimes_L \Z_{(p)}/F)$ and then mod out by $[p]$ in a second step.
Note that $W$ is a filtered colimit of projective generators (Remark~\ref{app:rw-projective}), so its derived tensor products are computed by ordinary tensor products.
The symmetric powers of $W$ split into sums of twists of $W$ (Remark~\ref{app:rw-projective}).
Consequently derived tensor products of these objects are also zero-truncated.
We find that the bar construction modelling $\dRW_{(p)} \hotimes_L \Z_{(p)}$ is a simplicial object with entries ordinary (underived) Hopf algebras, all of which are in fact sums of twists of $W$.
To determine the cofiber of $F$ on $\dRW_{(p)} \hotimes_L \Z_{(p)}$, we may take the cofiber of $F$ termwise in the bar construction.
By Lemma \ref{lemm:W-mod-F}, this turns every summand $W$ into $\Ga$.
It follows that the simplicial Hopf ring $\dRW_{(p)} \hotimes_L \Z_{(p)}/F$ is modeled by the simplicial object which is the \emph{underived} bar construction corresponding to 
\begin{equation*}
\begin{CD}
\Sym_{\F_p} \phi_p^* \Ga @>>> \F_p \\
  @V\Sym(0)VV        @VVV \\
\Sym_{\F_p} \Ga @>>> (\dRW_{(p)} \hotimes_L \Z_{(p)})/F.
\end{CD}
\end{equation*}
Using the equivalence between Hopf algebras with $F=0$ and vector spaces with Verschiebung (Proposition \ref{prop:indecomp-charp}), Dold--Kan identifies this simplicial object with $\Sym(\Ga) \hotimes \Lambda(\phi_p^* \Ga)$.
Here and below, $\Sym$ and $\Lambda$ without subscripts means over $\F_p[\F_p]$ (which only affects the degree zero part).
As usual, $\Sym(\Ga)$ is placed in chain degree zero, and $\Lambda^1(\Ga)$ in degree one.
The differential is zero (since the maps inducing the above square are), and so we have found that \[ \hH_*(\dRW_{(p)} \hotimes_L \Z_{(p)}/F) \wequi \Sym(\Ga) \hotimes \Lambda(\phi_p^* \Ga). \]
In particular, $[p]=0$ on these Hopf algebras, from which we deduce that \[ \hH_*(\dRW_{(p)} \hotimes_L \F_p/F) \wequi \Sym(\Ga) \hotimes \Lambda(\phi_p^* \Ga) \hotimes \Lambda(\epsilon), \] where $\epsilon$ is placed in chain degree one and internal degree zero.

\emph{Step 2. Determining the $E_1$-page.}
In what follows, almost all our manipulations involve Hopf algebras with $F=0$, and so we will mainly be thinking about indecomposables.
We write $\alpha_{(0)}, \alpha_{(1)}$ etc. for the generators of $\Sym^1(\Ga)$.
We will systematically need to deal with also $\phi_p^*\Sym^1(\Ga)$, $\phi_p^{2*}\Sym^1(\Ga)$, and so on.
As in Example \ref{ex:Ga-hotimes-Ga}, we will use the (surjective but not injective) Verschiebung maps to identify generators of the twists.
Thus e.g. $\phi_p^* \Ga$ has generators $\alpha_{(1)}, \alpha_{(2)}$, and so on.
We symbolically denote this correspondence as $\phi_p^* \Ga \wequi \Ga/\alpha_{(0)}$, and so on.
We denote the generators of $\Lambda^1(\phi_p^* \Ga)$ by $\beta_{(1)}, \beta_{(2)}$, and so on, and similarly for $\Lambda^1 \phi_p^{2*}\Ga$, etc.
Let us denote by $G^*$ the associated graded of the (homology of the) filtration by Frobenius on $\dRW_{(p)} \hotimes_L \F_p$.
Since $VF = [p] = FV$, Verschiebung induces a map $G^i \to G^{i+1}$ (not to be confused with the internal Verschiebung on each $G^i$!).
We will view $G^*$ as an $\F_p[t]$-module via this operator.
As such, we identify\NB{details?} \[ G^* \wequi \Sym(\Ga) \hotimes \Lambda(\Ga) \hotimes \Lambda(\epsilon) \hotimes \F_p[t] / (\beta_{(0)}, t\alpha_{(0)}, t\beta_{(1)}, t^2\alpha_{(1)}, \dots, t^i\alpha_{(i-1)},t^i \beta_{(i)}, \dots). \]
This is the $E_1$-page of our spectral sequence, converging strongly to $\hH_*(\dRW_{(p)} \hotimes_L \F_p)$ (at least on connected parts).
Its $d_1$-differential is induced by the filtration from the original differentials in the bar construction.
Using again that $Fa \hotimes b = F(a \hotimes Vb)$ one checks\NB{details} that the differential is determined by \[ d(\beta_{(i)}) = t\alpha_{(i)}, \quad d(\alpha_{(i)}) = 0, \quad d(\epsilon) = t. \]

\emph{Step 3. Studying the spectral sequence.}
Write \[ \tilde G^* = \Sym(\Ga) \hotimes \Lambda(\Ga) \hotimes \Lambda(\epsilon) \hotimes \F_p[t] \] with differential as above, so that there is an evident surjection $\tilde G^* \to G^*$.
We shall show at the end that $\tilde G^* \to G^*$ is a quasi-isomorphism.
Assuming this, write \[ \tilde G_0^* = \Sym(\Ga) \hotimes \Lambda(\Ga)\hotimes \F_p[t], \] also with the above differential.
Then we see that \[ \tilde G^* \cong \tilde G_0^* \otimes^L_{\F_p[t]} \F_p \wequi \tilde G_0^* \otimes_{\F_p[t]} \F_p, \] the latter equivalence since $\tilde G_0^*$ consists of free $\F_p[t]$-modules.
But \[ G_0^* \otimes_{\F_p[t]} \F_p \wequi \Sym(\Ga) \hotimes \Lambda(\Ga) \] placed entirely in $*=0$, so in particular with zero differential.
This proves what we want.

\emph{Step 4. Showing that $\tilde G^* \to G^*$ is a quasi-isomorphism.}
We work entirely on indecomposables.
Let $A_0 \subset \tilde G^*$ be the submodule spanned by $\beta_{(0)}$, so that \[ A_0 = \beta_{(0)} \wedge \Sym(\alpha_{(i)}, i \ge 0) \otimes \Lambda(\beta_{(j)}, j \ge 1) \otimes \Lambda(\epsilon) \otimes \F_p[t]. \]
For $\beta_{(0)} \wedge F \in A_0$ we have $d(\beta_{(0)} \wedge F) = t\alpha_{(0)} F \pm \beta_{(0)} \wedge dF$ which vanishes only if $F=0$.
Hence $d: A_0 \to dA_0$ is an isomorphism and $\tilde G^* \to \tilde G^*/(\beta_{(0)}, t\alpha_{(0)})$ is a quasi-isomorphism.
Similarly we can let $A_1 \subset \tilde G^*/(A_0, dA_0)$ be the submodule spanned by $t\beta_{(1)}$ and argue as before\NB{details?} to see that $\tilde G^*/(\beta_{(0)}, t\alpha_{(0)}) \to \tilde G^*/(\beta_{(0)}, t\alpha_{(1)}, t\beta_{(1)}, t^2\alpha_{(1)})$ is a quasi-isomorphism.
Continuing this way (using that quasi-isomorphism are stable under filtered colimits) we obtain the desired result.
\end{proof}

\subsection{A regular-sequence calculation}\label{app:rw-koszul}

% COPIED from wilson.tex, lines 2267--2302; source comments retained.
We now discuss an exactness property.
We discuss it purely as a piece of algebra.
Consider the Hopf ring \[ \scr K = \Lambda(\Ga) \hotimes \Sym(\Ga) \hotimes \F_p[\F_p[a_0, a_1, a_2, \dots]] \in \Alg(\Hopf_{\F_p}) \] (which is not commutative because of the exterior algebra part).
There are canonical ``coordinates'' \[ x: \Ga \wequi \Lambda^1 \Ga \to \scr K \quad\text{and}\quad y: \Ga \wequi \Sym^1 \Ga \to \scr K. \]
We shall promote $\scr K$ to a differential graded Hopf ring, so in particular, a chain complex in $\Hopf_{\F_p}$.
To this end, we declare that $x$ has degree $1$ and everything else has degree $0$, so that \[ \scr K_n = \Lambda^n(\Ga) \hotimes \Sym(\Ga) \hotimes \F_p[\F_p[a_0, a_1, a_2, \dots]] \in \Hopf_{\F_p}. \]
The differential $d: \scr K_* \to \scr K_{*-1}$ is induced by \[ d(x) = \sum_i a_i y^{p^i}. \]
In other words, we are looking at the Koszul complex of the ``element'' $d(x)$.

\begin{remark} \label{rmk:build-complex}
Some explanations are in order.
First note that $d(x)$ is well-defined, since $y: \Ga \to \Sym(\Ga) \hotimes \F_p[\F_p[a_0, a_1, a_2, \dots]]$ is pointed and $\Ga$ is ind-conilpotent (using ind-conilpotence of $\A^1$); it is a map of Hopf algebras by the Frobenius identity for Hopf-ring multiplication.
Now we define a map
\begin{gather*}
\tilde d_n: \Ga^{\htimes n} \longrightarrow
\Lambda(\Ga) \hotimes \Sym(\Ga) \hotimes \F_p[\F_p[a_0,a_1,\dots]],\\
\tilde d_n = \sum_i [(-1)^{i-1}]
 p_1\cdot\dots\cdot p_{i-1}\cdot (d_1\circ p_i)
 \cdot p_{i+1}\cdot\dots\cdot p_n.
\end{gather*}
Here $p_i: \Ga^{\htimes n} \to \Ga$ is projection to the $i$-th factor, $d_1$ is $d(x)$, and the multiplication uses the ring structure of $\scr K$.
Then $\tilde d_n$ is multilinear, so passes to $\Ga^{\hotimes n}$, and moreover satisfies ``$d(a \hotimes a) = 0$'', so passes to the exterior power (by definition of the exterior power).
Extending linearly over the remaining tensor factors, we obtain $d_n: \scr K_n \to \scr K_{n-1}$.
This is easily verified to be a differential (i.e. square to zero).
\end{remark}

\begin{lemma} \label{lemm:regular-sequence}
The complex $\scr K$ is acyclic.
In other words, for $n > 0$, the sequence of Hopf algebras \[ \scr K_{n+1} \xrightarrow{d} \scr K_n \xrightarrow{d} \scr K_{n-1} \in \Hopf_{\F_p} \] is exact.
\end{lemma}
\begin{proof}
The internal degree zero part of $\scr K$ is concentrated in chain degree zero, and so acyclic.
It thus suffices to consider the connected part of $\scr K$.
Since $F=0$ on $\Ga$, the same is true on $\Ga^{\hotimes n}$ and thus on $\Lambda(\Ga), \Sym(\Ga)$ (see Example \ref{ex:Ga-hotimes-Ga}).
In other words, $F=0$ on all of $\scr K^\cn$.
To check exactness, we may thus apply the symmetric monoidal functor of indecomposables (Proposition \ref{prop:indecomp-charp}).
Now \[ Q(\Sym(\Ga) \hotimes \F_p[\F_p[a_0, a_1, a_2, \dots]]) \wequi \F_p[\alpha_{(0)}, \alpha_{(1)}, \dots] \otimes \F_p[a_0, a_1, a_2, \dots] \] and $Q(\scr K)$ is the Koszul complex of the sequence $r_1, r_2, \dots$, where \[ r_j = \sum_{i} a_i F^i(\alpha_{(j)}). \]
Here $F^i(\alpha_{(j)})$ means the $i$-fold Hopf ring Frobenius applied to $\alpha_{(j)}$.
To determine $F(\alpha_{(j)})$, we must consider $\Delta_p(\alpha_{(j)}) \in \Ga^{\hotimes p}$ and its image in $Q(\Sym^p \Ga)$.
Using the formula for the coproduct on $\Ga$ (see Example \ref{ex:Ga-hotimes-Ga}) we see that $F(\alpha_{(j)})=\alpha_{(j-1)}^p$ and thus $F^i(\alpha_{(j)}) = \alpha_{(j-i)}^{p^i}$.
(In particular the sum is finite, as we knew it had to be.)
Regularity of this sequence is proved by Ravenel--Wilson \cite[Lemma 4.15(b)]{ravenel-wilson}.\NB{They do not have $a_0$, but this is irrelevant since our $r_0$ is just $a_0$...}
\end{proof}

\subsection{The comparison}\label{app:rw-comparison}

Put $I=\ker(L\to\Z_{(p)})$ and $J=\ker(L\to\F_p)=(p,I)$.

\begin{lemma}\label{lemm:coefficients}
Write $[p]_{\mathcal F}(T)=\sum_{m\ge1}c_mT^m\in L\fpsr{T}$. All $c_m$ lie in $J$, so $[p]_{\mathcal F}(T)\in J\fpsr{T}$, and modulo $J^2$ we have, in $(J/J^2)\fpsr{T}$,
\[ [p]_{\mathcal F}(T)=\sum_{i\ge0}a_iT^{p^i},\qquad a_0=\bar p. \]
The elements $a_i\in J/J^2$ are linearly independent and can be completed to a homogeneous basis. Consequently
\[\operatorname{gr}_J L\simeq \F_p[a_0,a_1,\ldots,b_1,b_2,\ldots].\]
\end{lemma}
\begin{proof}
This follows from standard formal group law theory; see, e.g., \cite[Appendix A2.1--A2.2]{ravenel1986complex}.
\end{proof}

In order to prove our main theorem, we filter $\dRW_{(p)}$ by powers of $J$ and study the resulting spectral sequence.
The coefficient exact sequences \[0\to J^n/J^{n+1}\to L/J^{n+1}\to L/J^n\to0\] give a Bockstein exact couple, and we now determine the first differential.
\begin{lemma}\label{lemm:bockstein}
Put $S=\dRW_{(p)}\hotimes_L \F_p$, and let $y$ and $x$ be the generators of $\hH_0(S)^1$ and $\hH_1(S)^1$ from Lemma~\ref{lemm:filter-derived-RW}. The first differential in the above exact couple is the derivation determined by
\[d_1(x)=\sum_{i\ge0}a_i y^{p^i},\qquad d_1(y)=0,\qquad d_1(a)=0\quad(a\in\operatorname{gr}_J L).\]
\end{lemma}
\begin{proof}
Let $\tilde y:\A^1\to \dRW_{(p)}^1$ be the canonical $[p]$-typical coordinate. The defining relation gives $\frob_p^{\mathcal F}(\tilde y)=0$. Applying $\ver_p^{\mathcal F}$ yields
\[ [p]_{\mathcal F}(\tilde y)=0\quad\text{in }\hH_0(\dRW_{(p)}). \]
Since every coefficient of the $p$-series belongs to $J$, the expression
\[r=\sum_m c_m\tilde y^m\in\hH_0(J\hotimes_L \dRW_{(p)})\]
is defined by conilpotence and maps to zero in $\hH_0(\dRW_{(p)})$. The coefficient fiber sequence $J\hotimes_L \dRW_{(p)}\to \dRW_{(p)}\to S$ thus gives a class $z\in\hH_1(S)^1$ with $\partial_J(z)=r$.

We identify $z$ by passing to the additive specialization $L\to\Z_{(p)}$. Write $Z=\dRW_{(p)}\hotimes_L\Z_{(p)}$. By Lemma~\ref{lemm:filter-derived-RW}(1), $Z^1\simeq\Ga$, with the image of $\tilde y$ the identity curve. The connecting homomorphism for multiplication by $p$ is therefore an isomorphism
\[ \partial_p:\hH_1(S)^1\xrightarrow{\simeq}\hH_0((p)\hotimes_{\Z_{(p)}}Z)^1,\]
and is normalized by $\partial_p(x)=p\otimes y$. Naturality of the connecting homomorphism identifies the image of $r$ with $p\otimes y$, since $[p]_{\mathcal F}(T)$ becomes $pT$. Thus $z=x$ and
\[\partial_J(x)=\sum_m c_m\tilde y^m.\]
Passing to $J/J^2$, Lemma~\ref{lemm:coefficients} gives the asserted formula for $d_1(x)$. The class $y$ lifts to $\dRW_{(p)}$, and the coefficient classes lift from $L$, so their differentials vanish. Multiplicativity of the exact couple determines the differential on all of its first page.
\end{proof}

\begin{lemma}\label{lemm:rw-positivity}
The normalized bar complex $\dRW_{(p)}$ is positive: $ (\dRW_{(p)}^r/\F_p)(s)=0$ for $s<r$. For a positive complex, its derived $I$-adic quotient tower ($I=\ker(L\to\Z_{(p)})$) is degreewise eventually constant. The connected part of a connective graded Hopf-algebra complex over $\F_p$ is $[p]$-complete.
\end{lemma}
\begin{proof}
For the first two assertions, positivity holds on $L$ and on $W$ placed in external degree one, and is preserved by tensor products, symmetric powers, and the bar realization. In any fixed bidegree, a sufficiently high product of elements of $I$ comes from external degree exceeding the internal degree; the same bound applies to a bar resolution of a derived $I$-quotient. For the last assertion one reduces to graded Hopf algebras and uses $[p]=FV=VF$; on the connected part, iterations of Verschiebung lower the positive internal degree until they vanish. Hence the localization at $[p]$ is zero and the connected part is $[p]$-complete. Compare \cite[Proposition 9.1]{goerss-dieudonne}.
\end{proof}

\begin{lemma}\label{lemm:joint-completeness}

The connected part $\dRW_{(p)}^{\cn}$ is derived $J$-complete: the canonical map
\[\dRW_{(p)}^{\cn}\longrightarrow\lim_n(\dRW_{(p)}\hotimes_L L/J^n)^{\cn}\]
is an equivalence.
\end{lemma}
\begin{proof}
The defining bar construction shows, by Lemma~\ref{lemm:rw-positivity}, that $\dRW_{(p)}$ is positive:
\[\big(\dRW_{(p)}^r/\F_p\big)(s)=0\qquad(s<r).\]
Put $K_n=(p^n,I^n)$. The inclusions
\[J^{2n}\subseteq K_n\subseteq J^n\]
show that the $K_n$-adic and $J$-adic quotient towers are pro-equivalent.

Choose homogeneous polynomial generators $t_j$ for $I$, all in strictly negative external degrees. The degree bound of Lemma~\ref{lemm:rw-positivity} applies uniformly to their products: in any fixed external and internal bidegree, the contribution of $I^n$ to a positive object vanishes for $n\gg0$. The same degree bound applies to derived $I$-quotients, obtained from the bar construction, and holds uniformly after reduction modulo $p^n$. Thus the tower
\[(\dRW_{(p)}/p^n)^{\cn}\longrightarrow ((\dRW_{(p)}/p^n)\hotimes_L L/I^n)^{\cn}\]
is degreewise eventually an equivalence as $n\to\infty$. Since $L/I^n$ is free over $\Z_{(p)}$, derived base change identifies the right hand side with $(\dRW_{(p)}\hotimes_L L/K_n)^{\cn}$. Consequently
\[\lim_n(\dRW_{(p)}\hotimes_L L/J^n)^{\cn}\simeq\lim_n(\dRW_{(p)}/p^n)^{\cn}.\]
The connected part of $\dRW_{(p)}$ is $[p]$-complete by Lemma~\ref{lemm:rw-positivity}, applied to the connected graded Hopf algebras underlying its external components (and passage to normalized simplicial complexes). This identifies the last limit with $\dRW_{(p)}^{\cn}$.
\end{proof}

We can now establish our main result.

\begin{proof}[Proof of Theorem \ref{thm:rw-bar-exact}.]
With the notation of Lemma~\ref{lemm:bockstein}, let
\[P_n=\dRW_{(p)}\hotimes_L L/J^{n+1},\qquad V_n=(J^n/J^{n+1})\hotimes_{\F_p} S.\]
There are fiber sequences $V_n\to P_n\to P_{n-1}$, and hence a multiplicative exact couple whose first page is
\[E_1^{s,t}= (J^s/J^{s+1})\otimes_{\F_p}\hH_t(S).\]
By Lemma~\ref{lemm:filter-derived-RW}, this is
\[\operatorname{gr}_J L\otimes_{\F_p}\Sym(\Ga)\hotimes\Lambda(\Ga).\]
The differential of Lemma~\ref{lemm:bockstein}, together with the choice of generators in Lemma~\ref{lemm:coefficients}, identifies this complex with $\scr K\otimes_{\F_p} \F_p[b_1,b_2,\ldots]$. Lemma~\ref{lemm:regular-sequence} therefore gives
\[E_2^{s,t}=0\qquad(t>0).\]

We explain why this implies vanishing of positive homology, without assuming convergence for the uncompleted ring. By Lemma~\ref{lemm:joint-completeness}, the connected part of the tower is complete and separated. Represent it by a complete filtered complex. If $z$ is a cycle in positive homological degree and filtration $s$, its leading term is a $d_1$-cycle. The vanishing above allows us to subtract a boundary to move $z$ into filtration $s+1$. Iterating gives a sequence of corrections in strictly increasing filtration whose sum converges, and whose differential is $z$. Hence $\hH_t(\dRW_{(p)}^{\cn})=0$ for $t>0$.

In internal degree zero the defining bar construction is just the coefficient ring $L$, and is already discrete. We conclude that $\hH_t(\dRW_{(p)})=0$ for $t>0$. Since $\hH_0(\dRW_{(p)})=R$, the augmentation is a homology equivalence.
\end{proof}
